\begin{tikzpicture}[x=1cm,y=1cm,every node/.append style={anchor=center,inner sep=.3333em}]
\draw[Rule,line width=.7pt] (0,0) circle(2.2);
\node[pNote] at(0,2.5){N};
\draw[-{Latex[length=2.4mm]},GeoA,line width=1pt] (0,0)--({2.2*sin(-10)},{2.2*cos(-10)}) node[pLabel,anchor=south east]{$350^\circ$};
\draw[-{Latex[length=2.4mm]},GeoA,line width=1pt] (0,0)--({2.2*sin(10)},{2.2*cos(10)}) node[pLabel,anchor=south west]{$10^\circ$};
\draw[-{Latex[length=2.4mm]},Classical,line width=1.6pt] (0,0)--(0,1.4);
\node[pLabel,text=Classical,anchor=west] at(.12,1.05){$\bar\theta=0^\circ$};
\draw[-{Latex[length=2.4mm]},GeoB,line width=1.1pt,dashed] (0,0)--(0,-2.2);
\node[pLabel,text=GeoB,anchor=west] at(.12,-1.6){media aritmetica $180^\circ$};
\fill[Ink] (0,0) circle(1.6pt);
\node[pNote,text width=6.4cm] at(0,-3.05){Risultante dei due vettori unitari: lunghezza $2\cos 10^\circ=1{,}97$, verso nord. Il modulo diviso per $n=2$ dà $\bar R=0{,}985$.};
\end{tikzpicture}
